% dynamics/realization_aware_dynamics.tex
\chapter{Realization-Aware Prediction and Comparison}
\label{chap:realization-aware-dynamics}

\section*{Which geometry enters the model?}

The intended Alignment geometry, its Metric Realization, the constructed
asset, physical state and observed-state estimate are related but not
interchangeable. A coordinate transformation changes representation; it does
not construct an asset. A survey polyline is an observation-bearing
representation; it is not physical truth.

\begin{principle}[Realization-aware dynamics principle]
\label{princ:realization-aware-dynamics}
A prediction shall identify which qualified realization or observed-state
estimate enters response model \(M\). The result remains a prediction of
\(M\), not an observation.
\end{principle}

\section{Typed comparison chain}

\begin{definition}[Realized runtime state]
\label{def:realized-runtime-state}
In this chapter, a realized runtime state is a response-model input derived
from an explicitly identified realization or observed-state estimate,
qualified by object identity, time, frame, method, provenance, applicability
and uncertainty. It is not a new canonical state class.
\end{definition}

For source \(q\in\{\mathrm{intent},\mathrm{realization},
\mathrm{observation}\}\),
\[
\widehat S_q
\xrightarrow[\;Q_M,I_M\;]{\mathcal D_M}
x_{M,q}^{\mathrm{pred}}(t).
\]
Only like-for-like outputs may be compared. A measured response
\(x^{\mathrm{obs}}(t)\) remains separately qualified:
\[
\Delta_M(t)=
\mathcal C\!\left(x_{M,q}^{\mathrm{pred}}(t),
x^{\mathrm{obs}}(t)\mid
\text{matched subject, time, frame and quantity}\right).
\]
\(\Delta_M\) is a comparison result, not a cause or decision.

\section{Realization filtering}

\begin{proposition}[Realization-filter principle]
\label{prop:realization-filter-principle}
Changing the qualified model input can change a prediction while preserving
Alignment identity. The change must not be attributed to construction,
degradation or measurement until evidence distinguishes those explanations.
\end{proposition}

For the CAD-curve/survey-polyline discrepancy, the CAD realization and the
polyline-derived estimate may each feed the same \(M\). Their output
difference is a sensitivity result. Causal attribution requires further
evidence.

\section{Calibration and uncertainty boundary}

A model parameter belongs to \(Q_M\). Calibration evidence supports selecting
or estimating such a parameter for a stated domain. Observation uncertainty
qualifies measured evidence; construction variability concerns the asset;
operational variability concerns use; a numerical residual concerns a
calculation; model-form uncertainty concerns the adequacy of \(M\). These
sources shall not be collapsed into one error term merely because all affect a
comparison.

Matching one observation can support a local calibration statement. It does
not establish universal model validity.

\begin{definition}[Realization-aware admissibility]
\label{def:realization-aware-admissibility}
Realization-aware admissibility is an evaluation result obtained by applying
an identified rule and purpose to qualified comparison evidence. It is not an
intrinsic property produced by \(\mathcal D_M\).
\end{definition}

\begin{definition}[Runtime envelope]
\label{def:runtime-envelope}
A runtime envelope is a declared set of limits from an identified source,
with units, applicability and validity. It is not supplied by the state model
itself.
\end{definition}

This yields inputs for optimization, but it does not choose an objective.
That responsibility is made explicit next.
